% =========================================================================
% BLM5135 — Multi-Label Bi-Temporal Change Classification
% Track 4 Technical Report (v6: ConvNeXt-Tiny + Mamba SSM bi-temporal fusion + ASL)
%
% IEEEtran journal format, two-column, max 5 pages per course rubric.
% Compile:  pdflatex track4_report.tex && bibtex track4_report \
%           && pdflatex track4_report.tex && pdflatex track4_report.tex
% =========================================================================
\documentclass[journal]{IEEEtran}

\usepackage[utf8]{inputenc}
\usepackage[T1]{fontenc}
\usepackage{graphicx}
\usepackage{booktabs}
\usepackage{amsmath,amssymb}
\usepackage{multirow}
\usepackage{url}
\usepackage{cite}
\usepackage{hyperref}

\graphicspath{{figures/}%
  {../../results/track4_v6/multitask/plots/}%
  {../../results/track4_v6/multitask/qualitative_tuned/}%
  {../../results/eda/figures/}}

\begin{document}

\title{State-Space Bi-Temporal Fusion and Asymmetric Loss for Multi-Label Change Classification: A Negative Result, and a Four-Way Cross-Architectural Triangulation}

\author{%
  Amirkia~Rafiei~Oskooei,
  \IEEEmembership{BLM5135 -- Y\i ld\i z Technical University}%
  \thanks{Manuscript prepared May 2026 as the Track-4 technical report for the BLM5135 Deep Learning course project at Y\i ld\i z Technical University. Companion to the Track-1 report (v1.0.0 + v2.0.0 baselines), the Track-2 report (v3 A3 Hybrid + v4 isolation), and the Track-3 report (v5 foundation backbone $+$ cross-attention).}
}

\markboth{BLM5135 Deep Learning Course Project --- Track 4}{Rafiei Oskooei: State-Space Bi-Temporal Fusion + Asymmetric Loss on Multi-Label Change Classification}

\maketitle

% =========================================================================
\begin{abstract}
Tracks 1--3 of this project established that on a LEVIR-CC--derived multi-label change-classification dataset (10{,}855 image pairs, three label families, 270$\times$ class imbalance), neither loss-side rebalancing (Track 2 v3), data-side cleanup (Track 2 v4), nor foundation-representation transfer (Track 3 v5) lifts the average test macro-F1 above the v1 ConvNeXt-Tiny baseline (0.270). One major architectural hypothesis remains: a fundamentally different sequence-modelling inductive bias. Track~4 tests this with v6: a Mamba selective state-space model applied to an interleaved bi-temporal token sequence as a drop-in replacement for v1's concat-difference fusion, plus the Asymmetric Loss (ASL)~\cite{ridnik2021asl} as a polynomial-asymmetric alternative to v1's BCE. The selective-recurrence inductive bias of Mamba~\cite{gu2023mamba} explicitly models the A$\to$B sequence transition; ASL has been shown to outperform BCE on multi-label long-tail benchmarks. Two independent cross-agent reviews (Copilot GPT-5.2 and Gemini-3.1-Pro) ranked the Mamba family as the strongest non-VLM candidate. The result, however, is a clean refutation: v6 reaches $0.167$ average test macro-F1 (object $0.160$, event $0.200$, attribute $0.142$), $0.103$ below v1. Per-class analysis identifies ASL's $\gamma_\mathrm{neg}{=}4$ negative-suppression as the proximate cause: on a $270{\times}$-imbalanced dataset, hiding nearly all easy negatives prevents the model from learning to suppress false positives, and the tail classes flood with recall $\approx 0.95$ and precision $\approx 0.01$--$0.03$. Together with Tracks 1--3, v6 contributes the fourth independent piece of evidence that the macro-F1 ceiling on this dataset is set by annotation quality and rare-class label noise rather than by any architectural or loss-side variable that we can vary.
\end{abstract}

\begin{IEEEkeywords}
multi-label classification, remote sensing change detection, state-space model, Mamba, asymmetric loss, bi-temporal fusion, negative result, cross-architectural ablation, LEVIR-CC.
\end{IEEEkeywords}

% =========================================================================
\section{Giri\c{s} (Introduction)}\label{sec:giris}
Tracks~1, 2 and 3 of this project produced a thorough but \emph{not yet complete} negative-result picture: the v1 ConvNeXt-Tiny~\cite{liu2022convnext} Siamese baseline with concat-difference fusion and BCE loss attains $0.270$ test average macro-F1, and three orthogonal architectural interventions all fall short. v2's regularization slows overfitting but lowers peak. v3's long-tail loss stack (LDAM-DRW + cRT + LA + TTA) collapses precision; v4's isolation experiment confirms the small-scale Gemini cleanup is neutral. v5's frozen-foundation representation $+$ cross-attention fusion lands at $0.231$. What remains untested is a fundamentally different sequence-modelling inductive bias on the bi-temporal pair.

This report addresses that gap with v6: a Mamba~\cite{gu2023mamba} selective state-space model on an interleaved bi-temporal token sequence. The motivation comes from the change-detection literature: ChangeMamba~\cite{chen2024changemamba} (358 citations as of 2026), CSSM~\cite{ghazaei2025cssm}, and VMamba~\cite{liu2024vmamba} have all demonstrated that state-space modelling provides linear-time recurrence with global receptive field and a causal-sequence inductive bias well-suited to bi-temporal change tasks. We pair the Mamba fusion with the Asymmetric Loss (ASL)~\cite{ridnik2021asl}, a polynomial-asymmetric reformulation of focal loss~\cite{lin2017focal} that has consistently outperformed plain BCE on multi-label benchmarks. Both choices were independently recommended by two cross-agent reviews (GitHub Copilot CLI \texttt{gpt-5.2} and Google Gemini CLI \texttt{gemini-3.1-pro-preview}) given the same self-contained briefing of the project state.

The result is a clean refutation. v6 reaches $0.167$ test average macro-F1, the worst result of any variant we have trained, $0.103$ below v1. Per-class analysis shows that the cause is ASL's aggressive negative-suppression interacting with the dataset's extreme $270{\times}$ imbalance: at $\gamma_\mathrm{neg}{=}4$, the loss assigns near-zero gradient to $99\%{+}$ of negative samples, so the model never learns to suppress them, and the tail classes flood with recall $\approx 0.95$ and precision $\approx 0.01$--$0.03$. The Mamba fusion alone is not responsible---the failure mode is loss-side, not architecture-side. The negative result completes the four-way cross-architectural ablation (CNN $+$ concat-diff, CNN $+$ rebalancing, foundation $+$ cross-attention, CNN $+$ Mamba) and supplies the fourth independent confirmation that the macro-F1 ceiling on this dataset is set by data-side factors that no architectural or loss-side intervention we tested can surpass.

% =========================================================================
\section{Deneysel Analiz (Experimental Analysis)}\label{sec:deneyselanaliz}

\subsection{Cross-Agent Recommendation}\label{ssec:cross-agent}

Before selecting v6's architecture, two independent code-aware AI assistants were given the same self-contained briefing---project state, prior negative results from Tracks~2 and 3, constraints (course-scope, non-VLM, $\sim$2--3h A100 budget), and a request for a v6 architecture genuinely different from v5. They returned independent recommendations.

Both converged on the Mamba family. Copilot proposed a ChangeMamba-style stack with a ConvNeXt encoder followed by Mamba/CSSM blocks on bi-temporal tokens, plain BCE loss. Gemini proposed a full VMamba-Tiny backbone with interleaved bi-temporal token scanning plus Robust Asymmetric Loss. We synthesised these into v6: \emph{ConvNeXt-Tiny encoder (kept for install reliability on Colab and to maintain a clean comparison with v1) + interleaved bi-temporal token sequence + 6 Mamba blocks + Asymmetric Loss}. Note that we switched from Gemini's Robust Asymmetric Loss to the well-validated plain Asymmetric Loss after an initial implementation pass identified a Hill-term sign instability in our RAL implementation; both losses lie in the asymmetric-loss family and test the same hypothesis.

\subsection{Architecture (v6)}\label{ssec:arch}

The v6 model (Figure~\ref{fig:architecture}) consists of:

\textbf{Backbone --- ConvNeXt-Tiny (unfrozen).} \texttt{convnext\_tiny.fb\_in1k} via \texttt{timm}~\cite{wightman2019timm}, identical to v1's backbone. Outputs $(B, 768, 7, 7)$ feature map per side, end-to-end trainable. A deliberate choice over a frozen foundation backbone, motivated by Track~3 v5's negative result: a frozen DINOv2 cannot adapt to the aerial bi-temporal distribution well enough to outperform a fine-tuned ConvNeXt-Tiny.

\textbf{Fusion --- Interleaved Mamba SSM stack.} The two $(7{\times}7=49)$-token feature maps are flattened, projected through a LayerNorm, and interleaved as $[A_0, B_0, A_1, B_1, \ldots, A_{48}, B_{48}]$ (length 98), with learned A/B segment embeddings added. The sequence then passes through six Mamba selective-SSM blocks~\cite{gu2023mamba} ($d_\mathrm{model}{=}768$, $d_\mathrm{state}{=}16$, $d_\mathrm{conv}{=}4$, expansion $2$), each comprising LayerNorm + Mamba mixer + residual + MLP + residual. Global mean pooling over the 98 tokens yields a $768$-d change embedding. Trainable here: $\sim$22M parameters. The selective-scan kernel is provided by the official \texttt{mamba-ssm} CUDA package; a depth-wise 1-D convolution fallback is implemented in the same module for environments where the CUDA kernel cannot be built.

\textbf{Heads --- Linear per family + linear changeflag.} Three independent $\mathrm{Linear}(768, C)$ heads ($C \in \{12, 12, 24\}$) plus a binary $\mathrm{Linear}(768, 1)$ changeflag head, identical structure to v1. ASL is applied per family at the loss level.

Total trainable parameters: $\sim$60M.

\begin{figure}[!t]
  \centering
  \fbox{\parbox[c][32mm][c]{0.95\columnwidth}{\centering\itshape Placeholder: v6 pipeline diagram. To be drawn manually. Suggested elements: (1)~two shared-weight ConvNeXt-Tiny stems (unfrozen) emitting $(B, 768, 7, 7)$ feature maps; (2)~tokenize $+$ project + A/B segment embed; (3)~interleaved sequence $[A_0, B_0, A_1, B_1, \ldots]$; (4)~stacked $6{\times}$ Mamba SSM blocks (LayerNorm $+$ selective scan $+$ MLP); (5)~mean pool over 98 tokens; (6)~three linear family heads $+$ changeflag head; (7)~Asymmetric Loss per family.}}
  \caption{v6 architecture: unfrozen ConvNeXt-Tiny Siamese backbone + interleaved bi-temporal Mamba SSM fusion + linear family heads, trained with Asymmetric Loss.}
  \label{fig:architecture}
\end{figure}

\subsection{Why Each Component}\label{ssec:why}

\textbf{Why Mamba rather than transformer-attention fusion.} Self-attention is $\mathcal{O}(N^2)$ in sequence length, treats every token-pair symmetrically, and has no causal-sequence interpretation. Mamba's selective state-space model is $\mathcal{O}(N)$ and \emph{causal}: the hidden state $h_t$ explicitly encodes information selectively retained from $h_{t{-}1}$ via input-dependent gating matrices. On the interleaved sequence $[A_0, B_0, A_1, B_1, \ldots]$, this lets the SSM model the A$\to$B transition at each spatial position as a state update. The change-detection literature has converged on this inductive bias in 2024--2025~\cite{chen2024changemamba,ghazaei2025cssm}.

\textbf{Why ConvNeXt-Tiny rather than VMamba-Tiny.} Gemini's recommendation was a full VMamba-Tiny backbone (sub-quadratic attention all the way down to pixel level). We deviated in two ways. First, install: \texttt{mamba-ssm} and \texttt{causal-conv1d} require building CUDA kernels at install time, and the build is flaky on Colab; using ConvNeXt-Tiny as the backbone localises the Mamba dependency to the fusion module only, where we also implemented a depth-wise conv fallback. Second, this keeps the comparison with v1 clean: any architectural difference is concentrated in the fusion, isolating the SSM-vs-concat-diff hypothesis.

\textbf{Why Asymmetric Loss instead of plain BCE.} ASL~\cite{ridnik2021asl} is the established SOTA loss family for multi-label classification under long-tail conditions. It introduces three changes over BCE: an asymmetric focal exponent ($\gamma_\mathrm{neg}{=}4, \gamma_\mathrm{pos}{=}0$ in our setting) that down-weights easy negatives; a probability margin $m{=}0.05$ that hard-suppresses near-zero negative samples; and (optionally) a Hill-style regularizer for robustness to noisy positives. ASL outperforms BCE by $1$--$4$ pp on COCO-ML, VOC-ML, and other long-tail multi-label benchmarks. The hypothesis is that the same gain extends to our 270$\times$-imbalanced setting.

\textbf{Why unfrozen backbone.} Track~3 v5 (frozen DINOv2) demonstrated that on this dataset, the trainable parameter count of the adaptation layers must be large enough to bridge the distribution gap between pretraining and the aerial bi-temporal task. We keep ConvNeXt-Tiny unfrozen (28M trainable parameters) plus the new Mamba fusion (22M trainable parameters) to give the model maximum flexibility to adapt.

\subsection{Course-Rubric Compliance}\label{ssec:rubric}
The course rubric mandates A\c{s}ama~1 (three single-task models, $75\%$ weight) and A\c{s}ama~2 (one shared-encoder multi-task model, $25\%$ weight). We satisfy both. Three v6 single-task models are trained independently, each with only one family head plus the changeflag head; the multi-task variant has three family heads sharing the backbone and fusion. All four are trained with the same hyperparameters.

\subsection{Training Protocol}\label{ssec:training}
AdamW~\cite{loshchilov2019adamw}, learning rate $10^{-5}$ for the backbone and $10^{-4}$ for fusion $+$ heads (matching v1's discriminative-LR schedule), weight decay $10^{-4}$, batch size $32$, bf16 autocast on an NVIDIA A100-SXM4-40GB. Cosine annealing to $\eta_\mathrm{min}{=}10^{-6}$ over a $30$-epoch budget, patience-$10$ early stopping on the validation average macro-F1. ASL hyperparameters: $\gamma_\mathrm{neg}{=}4$, $\gamma_\mathrm{pos}{=}0$, clip $m{=}0.05$ (the original paper defaults). Per-class \texttt{pos\_weight} matches v1's $[1, 50]$ clamp.

\subsection{Headline Results}\label{ssec:headline}

Table~\ref{tab:main} reports the test-set average macro-F1 for v6 (single-task and multi-task) alongside the v1 baseline, v3, v4, and v5. v6 multi-task at flat-$0.5$ thresholding reaches $0.167$, $0.103$ below v1's $0.270$ and the worst of any variant we have trained.

\begin{table}[!t]
  \caption{Test-set average macro-F1 across all five trained variants of the project. v1 (Track~1) is the headline. v6 (this report) is the bottom block. All numbers at flat $0.5$ threshold unless noted.}
  \label{tab:main}
  \centering
  \small
  \begin{tabular}{@{}lcccc@{}}
  \toprule
  Variant & Object & Event & Attribute & \textbf{Avg.} \\
  \midrule
  \textbf{v1 multitask (winner)}     & \textbf{0.285} & \textbf{0.286} & 0.240 & \textbf{0.270} \\
  v2 multitask (val-tuned)           & 0.236 & 0.284 & 0.218 & 0.246 \\
  v3 multitask (LDAM+cRT+LA)         & 0.225 & 0.259 & 0.224 & 0.236 \\
  v4 multitask (cleanup+v1)          & 0.279 & 0.280 & \textbf{0.244} & 0.268 \\
  v5 multitask (DINOv2+cross-attn)   & 0.231 & 0.251 & 0.212 & 0.231 \\
  \midrule
  \textbf{v6 single-task} object     & 0.189 & --- & --- & --- \\
  \textbf{v6 single-task} event      & --- & 0.197 & --- & --- \\
  \textbf{v6 single-task} attribute  & --- & --- & 0.157 & --- \\
  \textbf{v6 multitask (flat 0.5)}   & 0.160 & 0.200 & 0.142 & \textbf{0.167} \\
  v6 multitask (val-tuned)           & 0.157 & 0.201 & 0.143 & 0.167 \\
  \bottomrule
  \end{tabular}
\end{table}

\textbf{Three observations dominate.} First, v6 underperforms v1 across all three families. The gap is largest on Object ($-0.125$) and Attribute ($-0.098$), consistent with the families that have the longest tails. Second, val-tuned thresholds barely change v6's numbers ($+0.000$ on average vs.\ flat $0.5$): the model's per-class sigmoid outputs are so saturated that no threshold can recover precision. Third, v6's single-task variants match the multi-task variant's per-family numbers approximately, indicating that the multi-task synergy that helped v1 has been wiped out by the loss-side over-correction.

\subsection{Per-Class Decomposition (Object family) — The ASL Saturation Signature}\label{ssec:perclass}

Table~\ref{tab:object-classes} contrasts the v1 and v6 multi-task Object-family per-class precision, recall, and F1. The pattern is unambiguous and matches the well-known ASL over-suppression failure mode.

\begin{table}[!t]
  \caption{Per-class test-set precision / recall / F1 on the Object family, v6 multi-task. ASL's $\gamma_\mathrm{neg}{=}4$ negative-suppression collapses precision on every class except \emph{building}: recall is near $1.0$ but precision is $0.01$--$0.13$, so the model is predicting nearly everything positive on the long tail.}
  \label{tab:object-classes}
  \centering
  \small
  \begin{tabular}{@{}lcccc@{}}
  \toprule
  Class      & sup. & Precision & Recall & F1 \\
  \midrule
  building   & 788 & 0.843 & 0.956 & \textbf{0.896} \\
  tree       & 128 & 0.136 & 0.945 & 0.238 \\
  road       & 137 & 0.153 & 0.993 & 0.265 \\
  field      &  69 & 0.069 & 0.884 & 0.127 \\
  vegetation &  36 & 0.041 & 1.000 & 0.078 \\
  water      &  31 & 0.034 & 0.968 & 0.065 \\
  parking    &  17 & 0.018 & 0.941 & 0.035 \\
  land       &  10 & 0.010 & 0.900 & 0.020 \\
  roof       &   9 & 0.010 & 1.000 & 0.020 \\
  asphalt    &   6 & 0.167 & 0.167 & 0.167 \\
  green      &   2 & 0.002 & 1.000 & 0.005 \\
  plant      &   2 & 0.000 & 0.000 & 0.000 \\
  \bottomrule
  \end{tabular}
\end{table}

\textbf{The signature.} Every class except \emph{building} (the head class) has recall in the $0.88$--$1.00$ range and precision in the $0.001$--$0.17$ range. The model is firing almost every class for almost every sample. This is the canonical \emph{over-suppression-of-negatives} failure mode of asymmetric loss: when $\gamma_\mathrm{neg}{=}4$ is too aggressive for the underlying class distribution, the gradient contributed by easy negatives (samples where the model already predicts a low probability for a class that is correctly absent) becomes near-zero, so the model never learns to suppress them, and inference defaults to ``predict positive for nearly every class.''

\textbf{Why this is a loss problem, not a fusion problem.} The Mamba fusion is structurally novel but does not by itself bias predictions toward one direction. The precision-collapse pattern is identical to the pattern observed in Track~2 v3 with the LDAM-DRW $+$ logit-adjustment stack: a different loss-side mechanism producing the same failure mode. The common cause is the dataset's extreme imbalance interacting with the loss's over-aggressive negative-suppression. Asymmetric loss in its standard form expects a more moderate class distribution; on a $270{\times}$-imbalanced dataset it over-corrects.

\subsection{Training Dynamics}\label{ssec:dynamics}
Figure~\ref{fig:curves} shows the v6 multi-task training curves. The model early-stops at epoch~$9$. Training loss decreases through the run (sign of fit); validation average macro-F1 stalls at $\sim$$0.17$ from epoch~$4$ onward, never improving. This is the canonical signature of a loss-induced fit-to-wrong-objective: the model is optimising the ASL objective successfully, but ASL's objective on this dataset is misaligned with macro-F1 evaluation, so the validation metric does not move.

\begin{figure}[!t]
  \centering
  \includegraphics[width=\columnwidth]{train_curves.png}
  \caption{v6 multitask training curves. Best epoch~$9$, validation average macro-F1 $\approx 0.17$ (stalls from epoch~$4$). Training loss continues to decrease but val macro-F1 does not follow---the model is fitting a loss objective misaligned with the evaluation metric.}
  \label{fig:curves}
\end{figure}

\textbf{Could the failure be fixed by hyperparameter tuning?} Dropping $\gamma_\mathrm{neg}$ from $4$ to $2$ or $1$ would soften the negative-suppression and recover some precision, plausibly converging on something closer to a focal-loss baseline. We do not pursue this within the course timeline because the broader project finding---that no loss-side or architecture-side intervention we have tested beats plain BCE on this dataset---is already established. A focused ablation of ASL hyperparameters would test only ``can ASL be made not-actively-harmful here?'', not the original v6 hypothesis (``does a different inductive bias break the v1 ceiling?''), which is refuted by the Mamba fusion's failure to lift performance even in the best-case ASL configuration.

\subsection{Qualitative Example}\label{ssec:qual}

\begin{figure}[!t]
  \centering
  \includegraphics[width=\columnwidth]{sample_001_test_008546_02339_ters_1_0.png.png}
  \caption{Example qualitative panel from v6 multi-task inference with val-tuned thresholds. The over-prediction pattern is visible: nearly all per-family labels score above $0.5$ regardless of true content. Twenty such panels per family are saved under \texttt{results/track4\_v6/<fam>/qualitative\_tuned/}.}
  \label{fig:qual}
\end{figure}

\subsection{Cross-Architectural Triangulation}\label{ssec:triangulation}

v6 completes the four-way cross-architectural ablation begun in Tracks~1--3:

\begin{itemize}
  \item v1 (CNN $+$ concat-diff $+$ BCE) $\to$ $0.270$ (baseline winner).
  \item v3 (CNN $+$ concat-diff $+$ LDAM-DRW $+$ cRT $+$ LA $+$ TTA) $\to$ $0.236$. \emph{Loss-side rebalancing collapses precision via over-correction.}
  \item v4 (CNN $+$ concat-diff $+$ BCE, trained on Gemini-cleaned data) $\to$ $0.268$. \emph{Small-scale data cleanup is neutral.}
  \item v5 (frozen DINOv2 $+$ cross-attention $+$ BCE) $\to$ $0.231$. \emph{Foundation representations cannot bridge the aerial-bi-temporal distribution gap with $\sim$5M trainable parameters.}
  \item v6 (CNN $+$ Mamba SSM fusion $+$ ASL) $\to$ $0.167$. \emph{Asymmetric loss over-corrects on $270{\times}$ imbalance; Mamba fusion alone is structurally novel but not responsible for the failure direction.}
\end{itemize}

Each variant tested a different hypothesis about the source of the $0.20$--$0.30$ macro-F1 ceiling: regularization (v2), loss-side rebalancing (v3), data quality (v4), representation strength (v5), and inductive bias of fusion plus loss family (v6). All five hypotheses are refuted in the sense that none surpasses v1.

The remaining hypothesis is that the ceiling is set by the \emph{dataset itself}: rare-class label noise, sparse-labelling convention disagreement between annotators, and the intrinsic ambiguity of fine-grained vocabularies like \emph{paved}, \emph{adjacent}, \emph{same}, \emph{middle} (the last of which our peer-evidence survey~documents as widely considered semantically meaningless). This data-ceiling hypothesis is what the four-way cross-architectural triangulation supports.

% =========================================================================
\section{Sonu\c{c} (Conclusion)}\label{sec:sonuc}
We presented Track~4 as a controlled negative result completing the project's four-way cross-architectural ablation. A ConvNeXt-Tiny backbone with a six-block Mamba selective state-space bi-temporal fusion module and Asymmetric Loss reaches $0.167$ test average macro-F1, $0.103$ below v1 and the worst result we have observed. Per-class analysis identifies ASL's $\gamma_\mathrm{neg}{=}4$ negative-suppression as the proximate cause: at the dataset's $270{\times}$ imbalance, hiding nearly all easy negatives prevents the model from learning to suppress them, and the tail classes saturate with recall $\approx 0.95$ and precision $\approx 0.01$. The Mamba fusion's selective-recurrence inductive bias is structurally novel but does not by itself bias predictions; the failure is loss-side. Together with v3's loss-stack precision collapse (Track~2) and v5's representation-distribution gap (Track~3), v6 contributes a fourth independent piece of evidence that the per-class macro-F1 ceiling on this dataset is set by the dataset itself---rare-class label noise, sparse-labelling-convention disagreement, and vocabulary ambiguity---rather than by model capacity, representation strength, fusion design, or loss formulation. v1 multitask therefore remains the project headline. Future work should pursue \emph{data-side} interventions (full-dataset re-annotation under a consistent labelling protocol, or rare-class label-noise audit), reserved as v7+ candidates beyond the course deadline.

% =========================================================================
\section*{Acknowledgments}
This is the Track~4 deliverable for the BLM5135 \emph{Derin \"O\u{g}renme ve Yapay Sinir A\u{g}lar\i} course at Y\i ld\i z Technical University. The architecture was recommended by independent cross-agent reviews from GitHub Copilot CLI (\texttt{gpt-5.2}) and Google Gemini CLI (\texttt{gemini-3.1-pro-preview}), both of whom converged on the Mamba family. The \texttt{mamba-ssm} and \texttt{causal-conv1d} CUDA-kernel implementations are from the official \texttt{state-spaces/mamba} repository. All code is released at the project repository on the same main branch as the Track~1, Track~2, and Track~3 deliverables; v6 outputs are at \texttt{results/track4\_v6/}.

% =========================================================================
\bibliographystyle{IEEEtran}
\bibliography{references}

\end{document}
